% =============================================================
%  § 1.  Элементы теории чисел
% =============================================================

\section{Элементы теории чисел}
\label{sec:numtheory}

\subsection{Делимость, остатки и сравнения по модулю}

Целые числа --- самые знакомые объекты в математике, и кажется, что про них трудно сказать что-то новое. Однако именно из вопросов о \emph{делимости} целых чисел вырос целый раздел математики --- \term{теория чисел}, --- и именно её результаты сейчас работают внутри каждого мессенджера и банковской карты.

\begin{definition}{}{def:p1-1}
Говорят, что целое число $a$ \term{делится} на целое число $b \neq 0$ (или что $b$ \term{делит} $a$), если существует такое целое число $q$, что $a = b\cdot q$. Пишут $b \mid a$.
\end{definition}

Если $a$ не делится на $b$, то всё равно можно записать
\[
a = b\cdot q + r, \qquad 0 \le r < |b|.
\]
Это знакомое со школы \term{деление с остатком}. Число $q$ называется \emph{неполным частным}, а $r$ --- \term{остатком от деления} $a$ на $b$. Остаток обозначают также $a \bmod b$ (читается «$a$ по модулю $b$»).

\begin{example}{}{ex:p1-1}
$17 = 5\cdot 3 + 2$, поэтому $17 \bmod 5 = 2$.\\
$-17 = 5\cdot(-4) + 3$, поэтому $(-17)\bmod 5 = 3$ (остаток --- неотрицательное число!).
\end{example}

Простые числа --- «атомы» делимости.

\begin{definition}{}{def:p1-2}
Натуральное число $p > 1$ называется \term{простым}, если оно делится только на единицу и на само себя. Натуральное $n > 1$, не являющееся простым, называется \term{составным}.
\end{definition}

Первые простые числа: $2, 3, 5, 7, 11, 13, 17, 19, 23, 29, \ldots$. Простых чисел бесконечно много --- это доказал ещё Евклид (доказательство мы повторим в п.~\ref{subsec:euclid}).

Основная теорема арифметики говорит, что любое натуральное число $n > 1$ \emph{единственным} (с точностью до порядка) способом раскладывается в произведение простых:
\[
n = p_1^{\alpha_1} \cdot p_2^{\alpha_2} \cdots p_k^{\alpha_k}.
\]

\paragraph*{Сравнения по модулю.} В каждой задаче, где важна только «остаточная» информация (например, день недели в зависимости от номера дня в году), удобно вместо чисел рассматривать их остатки от деления на некоторое фиксированное число.

\begin{definition}{}{def:p1-3}
Пусть $m \ge 2$ --- натуральное число (\term{модуль}). Целые числа $a$ и $b$ называются \term{сравнимыми по модулю $m$}, если их разность $a - b$ делится на $m$. Пишут
\[
a \equiv b \modn{m}.
\]
\end{definition}

\noindent
Равносильное определение: $a \equiv b \modn{m}$ тогда и только тогда, когда $a$ и $b$ дают одинаковые остатки при делении на~$m$.

\begin{example}{}{ex:p1-2}
$23 \equiv 9 \modn 7$, потому что $23 - 9 = 14 = 2\cdot 7$.\\ В обоих случаях остаток от деления на 7 равен 2.
\end{example}

Сравнения --- очень удобный язык, потому что они \emph{ведут себя как обычные равенства}: их можно складывать, вычитать, умножать и возводить в степень.

\begin{theorem}{свойства сравнений}{thm:p1-1}
Если $a \equiv b \modn m$ и $c \equiv d \modn m$, то:
\begin{enumerate}
  \item $a + c \equiv b + d \modn m$;
  \item $a - c \equiv b - d \modn m$;
  \item $a \cdot c \equiv b \cdot d \modn m$;
  \item $a^n \equiv b^n \modn m$ для любого натурального $n$.
\end{enumerate}
\end{theorem}

\begin{proof}[Идея доказательства]
По определению $a - b = m\cdot k_1$ и $c - d = m\cdot k_2$. Тогда $(a + c) - (b + d) = m(k_1 + k_2)$ делится на $m$. Для умножения: $ac - bd = ac - bc + bc - bd = c(a-b) + b(c-d)$, что тоже делится на $m$. Возведение в степень --- следствие умножения. \qedhere
\end{proof}

С \emph{делением} ситуация сложнее. Нельзя «просто так» поделить обе части сравнения --- например, $6 \equiv 0 \modn 6$, но «сократив на 2», получим $3 \equiv 0 \modn 6$, что неверно. Когда деление законно --- мы разберём чуть позже, после знакомства с алгоритмом Евклида.

\begin{interestbox}
Сравнения встречаются повсюду в жизни. Часы --- арифметика по модулю 12 (или 24): если сейчас 10 часов, то через 5 часов будет $10 + 5 = 15 \equiv 3 \modn{12}$. День недели --- арифметика по модулю 7. Контрольная цифра ИНН, номер банковской карты, штрихкод книги --- все они вычисляются как сумма цифр (с разными весами) по некоторому модулю.
\end{interestbox}

% -----------------------------------------------------------------
\subsection{Алгоритм Евклида}
\label{subsec:euclid}

Одна из главных операций в теории чисел --- нахождение \term{наибольшего общего делителя} (НОД) двух чисел.

\begin{definition}{}{def:p1-4}
Наибольший общий делитель чисел $a$ и $b$ (не равных одновременно нулю) --- это наибольшее натуральное число, на которое делятся одновременно и $a$, и $b$. Обозначается $\NOD(a,b)$.

Если $\NOD(a,b) = 1$, то числа $a$ и $b$ называются \term{взаимно простыми}.
\end{definition}

Как найти $\NOD(a,b)$? Можно, например, разложить оба числа на простые множители и взять общие. Но при больших числах разложение на множители --- очень тяжёлая задача (об этом мы поговорим в конце параграфа). К счастью, ещё в III веке до н. э. Евклид предложил способ, который работает несравнимо быстрее.

\paragraph*{Геометрическая идея.} Представим себе прямоугольник со сторонами $a$ и $b$, причём $a > b$. Будем «выкладывать» в него квадраты со стороной $b$, пока это возможно. Когда останется полоска со сторонами $b$ и $r = a \bmod b$, повторим процедуру для неё, и так далее. Когда выкладывание закончится точным квадратом --- сторона этого последнего квадрата и есть $\NOD(a,b)$ (рис.~\ref{fig:euclid-geom}).

\begin{figure}[H]
\centering
\begin{tikzpicture}[scale=0.55, every node/.style={font=\small}]
  % Параметры для прямоугольника 48 x 30
  % 48 = 1*30 + 18 (квадрат 30x30, остаток 30x18)
  % 30 = 1*18 + 12 (квадрат 18x18, остаток 18x12)
  % 18 = 1*12 + 6  (квадрат 12x12, остаток 12x6)
  % 12 = 2*6 + 0   (два квадрата 6x6) -> НОД = 6
  \def\W{16} % 48/3
  \def\H{10} % 30/3

  % Основной прямоугольник 48x30 -> 16x10
  \draw[thick, prosvBlue] (0,0) rectangle (\W,\H);

  % Квадрат 30x30 -> 10x10 (слева)
  \fill[prosvLight, opacity=0.7] (0,0) rectangle (10,10);
  \draw[thick, prosvBlue] (10,0) -- (10,10);
  \node at (5,5) {\large $30\times 30$};

  % Остаток 18x30 -> 6x10 (справа). Делим: квадрат 18x18 -> 6x6 сверху
  \fill[prosvCream, opacity=0.8] (10,4) rectangle (16,10);
  \draw[thick, prosvBlue] (10,4) -- (16,4);
  \node at (13,7) {$18\times 18$};

  % Остаток 18x12 -> 6x4 снизу. Квадрат 12x12 -> 4x4 слева
  \fill[prosvLight!60!white, opacity=0.7] (10,0) rectangle (14,4);
  \draw[thick, prosvBlue] (14,0) -- (14,4);
  \node at (12,2) {$12\!\times\!12$};

  % Остаток 6x4 -> 2x4 справа. Делим: два квадрата 6x6 = 2x2
  \fill[prosvGold!30!white, opacity=0.8] (14,0) rectangle (16,2);
  \draw[thick, prosvBlue] (14,2) -- (16,2);
  \draw[thick, prosvBlue] (15,0) -- (15,2);
  \fill[prosvGold!30!white, opacity=0.8] (14,2) rectangle (16,4);
  \draw[thick, prosvBlue] (15,2) -- (15,4); 

  \node[font=\scriptsize] at (15,1) {6};
  \node[font=\scriptsize] at (15,3) {6};
  
  % Подписи сторон
  \draw[<->,>=stealth, prosvRed] (0,-0.6) -- (\W,-0.6);
  \node[below, prosvRed] at (\W/2,-0.7) {$a = 48$};
  
  \draw[<->,>=stealth, prosvRed] (-0.6,0) -- (-0.6,\H);
  \node[left, prosvRed] at (-0.7,\H/2) {$b = 30$};

  \draw[<->,>=stealth, prosvGreen] (16.4,0) -- (16.4,2);
  \node[right, prosvGreen] at (16.5,1) {$\NOD = 6$};
\end{tikzpicture}
\caption{Геометрическая иллюстрация алгоритма Евклида для пары $(48,\,30)$. Прямоугольник заполняется квадратами: $30{\times}30$, затем $18{\times}18$, затем $12{\times}12$ и, наконец, два квадрата $6{\times}6$. Сторона последнего квадрата равна $\NOD(48,30) = 6$.}
\label{fig:euclid-geom}
\end{figure}

Численно это соответствует последовательным делениям с остатком:
\begin{align*}
48 &= 1\cdot 30 + 18,\\
30 &= 1\cdot 18 + 12,\\
18 &= 1\cdot 12 + 6,\\
12 &= 2\cdot 6 + 0 \;\;\Longrightarrow\;\; \NOD(48,30) = 6.
\end{align*}

\noindent
\emph{Почему это работает?} Ключевое наблюдение: если $a = bq + r$, то $\NOD(a, b) = \NOD(b, r)$. Действительно, любое число, делящее $a$ и $b$, делит и $r = a - bq$; обратно, любое число, делящее $b$ и $r$, делит и $a = bq + r$. Значит, у пар $(a,b)$ и $(b,r)$ один и тот же набор общих делителей --- и, следовательно, один и тот же $\NOD$. На каждом шаге второй аргумент уменьшается; раньше или позже он станет нулём, и тогда первое число --- искомый НОД.

\begin{algorithmbox}[title={Алгоритм: нахождение НОД (Евклид)}]
\textbf{Вход:} натуральные числа $a, b$, $a \ge b > 0$.\\
\textbf{Выход:} $\NOD(a,b)$.
\begin{enumerate}
  \item Если $b = 0$, вернуть $a$.
  \item Иначе вычислить $r = a \bmod b$ и повторить с парой $(b, r)$.
\end{enumerate}
\end{algorithmbox}

На псевдокоде:
\begin{center}
\begin{minipage}{0.7\linewidth}
\ttfamily\small
\noindent
function gcd(a, b):\\
\quad while b $\ne$ 0:\\
\qquad a, b := b, a mod b\\
\quad return a
\end{minipage}
\end{center}

\paragraph*{Скорость работы.} Сколько шагов делает алгоритм? Можно доказать (теорема Ламе, 1844 г.): число делений в алгоритме Евклида для $\NOD(a,b)$ не превосходит примерно $5\log_{10}\min(a,b)$. Это значит, что для чисел в \emph{сотни цифр} (а такие числа реально используются в криптографии) НОД находится за несколько сотен делений --- доли секунды.

% -----------------------------------------------------------------
\subsection{Расширенный алгоритм Евклида}

Поразительный факт: Евклид нашёл не только $\NOD$, но и способ \emph{представить} его в виде «целочисленной линейной комбинации» исходных чисел.

\begin{theorem}{соотношение Безу}{thm:p1-2}
\label{thm:bezout}
Для любых натуральных $a, b$ существуют целые числа $x, y$ (\emph{не обязательно положительные}), такие что
\[
a\cdot x + b\cdot y = \NOD(a,b).
\]
\end{theorem}

Эти $x, y$ называются \term{коэффициентами Безу}. Найти их позволяет \term{расширенный алгоритм Евклида}: при делениях с остатком мы заодно «отслеживаем», как выражается каждый текущий остаток через исходные $a$ и $b$.

\paragraph*{Пример: те же 48 и 30.} Прокрутим цепочку делений «снизу вверх», выражая 6 (наш НОД) через всё бо\'{}льшие числа.
\begin{align*}
6 &= 18 - 1\cdot 12 && \text{(из третьего деления)}\\
  &= 18 - 1\cdot (30 - 1\cdot 18) = 2\cdot 18 - 1\cdot 30 && \text{(подставили } 12 = 30 - 18 \text{)}\\
  &= 2\cdot (48 - 1\cdot 30) - 1\cdot 30 = 2\cdot 48 - 3\cdot 30 && \text{(подставили } 18 = 48 - 30 \text{)}
\end{align*}
Получили: $\boxed{2\cdot 48 + (-3)\cdot 30 = 6}$. Коэффициенты Безу: $x = 2$, $y = -3$.

\begin{example}{}{ex:p1-3}
Проверим: $2\cdot 48 - 3\cdot 30 = 96 - 90 = 6 = \NOD(48,30)$. Совпадает.
\end{example}

Удобно вести расчёт в таблице:
\begin{center}
\begin{tabular}{|c|c|c|c|c|}
\hline
шаг & $r$ & $q$ & $x$ & $y$ \\\hline
0 & 48 & --- & 1 & 0\\
1 & 30 & --- & 0 & 1\\
2 & 18 & 1 & 1 & $-1$\\
3 & 12 & 1 & $-1$ & 2\\
4 & 6  & 1 & 2 & $-3$\\
5 & 0  & 2 & --- & ---\\\hline
\end{tabular}
\end{center}
В каждой строке выполняются равенства $r = 48 x + 30 y$. На предпоследней строке (где $r = \NOD = 6$) и считываются коэффициенты.

% -----------------------------------------------------------------
\subsection{Решение линейных сравнений}

Применим то, что получили, к решению уравнений в модулярной арифметике.

\paragraph*{Обратный элемент по модулю.} В обычной арифметике число, обратное к $a$, --- это $1/a$: при умножении даёт единицу. По модулю $m$ \term{обратным к $a$} называется такое целое $a^{-1}$, что
\[
a \cdot a^{-1} \equiv 1 \modn m.
\]
Например, по модулю 7: $3\cdot 5 = 15 = 2\cdot 7 + 1 \equiv 1 \modn 7$, значит, $3^{-1} \equiv 5 \modn 7$.

\begin{theorem}{}{thm:p1-5}
Обратный элемент к $a$ по модулю $m$ \emph{существует} тогда и только тогда, когда $\NOD(a, m) = 1$. При этом он находится с помощью расширенного алгоритма Евклида.
\end{theorem}

\begin{proof}[Идея]
Если $\NOD(a,m) = 1$, то по теореме~\ref{thm:bezout} существуют целые $x, y$, такие что $ax + my = 1$. Это сравнение по модулю $m$ принимает вид $ax \equiv 1 \modn m$, то есть $x \equiv a^{-1} \modn m$.

Обратно, если $a\cdot a^{-1} \equiv 1 \modn m$, то $a \cdot a^{-1} - 1 = m\cdot k$ для некоторого $k$, откуда любой общий делитель $a$ и $m$ должен делить $1$, --- значит, $\NOD(a,m) = 1$.
\end{proof}

\begin{example}{}{ex:p1-4}
Найдём $11^{-1} \modn{26}$. Применяем расширенный алгоритм Евклида к $(26,11)$:
\begin{align*}
26 &= 2\cdot 11 + 4,\\
11 &= 2\cdot 4 + 3,\\
4 &= 1\cdot 3 + 1,\\
3 &= 3\cdot 1 + 0.
\end{align*}
$\NOD = 1$. Раскручиваем:
\begin{align*}
1 &= 4 - 1\cdot 3 = 4 - (11 - 2\cdot 4) = 3\cdot 4 - 11\\
&= 3\cdot (26 - 2\cdot 11) - 11 = 3\cdot 26 - 7\cdot 11.
\end{align*}
Значит, $-7\cdot 11 \equiv 1 \modn{26}$, то есть $11^{-1} \equiv -7 \equiv 19 \modn{26}$. Проверка: $11\cdot 19 = 209 = 8\cdot 26 + 1$.
\end{example}

\paragraph*{Линейные сравнения.} Уравнение вида
\[
a\cdot x \equiv b \modn m
\]
называется \term{линейным сравнением}. Если $\NOD(a,m) = 1$, его решение находится «домножением на обратный»: $x \equiv a^{-1}\cdot b \modn m$.

% -----------------------------------------------------------------
\subsection{Быстрое возведение в степень}
\label{subsec:fastpow}

В криптографии нужно уметь вычислять выражения вида $a^n \bmod m$ для огромных $n$ (порядка $2^{1000}$). Если в лоб умножать $a$ на себя $n$ раз --- это никогда не закончится: даже за миллиард умножений в секунду на это ушло бы больше времени, чем существует Вселенная.

Спасает идея \term{«разделяй и властвуй»}, известная также как \term{бинарное возведение в степень}. Идея проста: разделим показатель пополам:
\[
a^n = \begin{cases}
\bigl(a^{n/2}\bigr)^2, & n\text{ чётно},\\[2pt]
a\cdot \bigl(a^{(n-1)/2}\bigr)^2, & n\text{ нечётно}.
\end{cases}
\]
Каждое такое деление пополам уменьшает показатель примерно вдвое, поэтому всего нужно около $\log_2 n$ шагов вместо~$n$ (рис.~\ref{fig:binpow}).

\begin{figure}[H]
\centering
\begin{tikzpicture}[
  level distance=1.0cm,
  level 1/.style={sibling distance=10cm},
  level 2/.style={sibling distance=5cm},
  level 3/.style={sibling distance=2.5cm},
  every node/.style={
    rounded corners=2pt,
    draw=prosvBlue,
    fill=prosvLight,
    font=\small,
    inner sep=3pt
  },
  edge from parent/.style={->, >=stealth, draw=prosvBlue, thick}
]
\node {$a^{13}$}
  child {node {$a^{13} = a\cdot a^{12}$}
    child {node {$a^{12} = (a^6)^2$}
      child {node {$a^6 = (a^3)^2$}
        child {node {$a^3 = a\cdot a^2$}}
      }
    }
  };
\end{tikzpicture}
\caption{Дерево бинарного возведения $a^{13}$: показатель уменьшается вдвое (с поправкой на чётность) на каждом шаге. Всего около $\log_2 13 \approx 4$ шагов.}
\label{fig:binpow}
\end{figure}

\noindent
Каждый «спуск» делит показатель примерно на 2, поэтому всего получается около $\log_2 n$ шагов. Для $n = 2^{1000}$ это \emph{тысяча} умножений вместо $2^{1000}$.

На псевдокоде:
\begin{center}
\begin{minipage}{0.78\linewidth}
\ttfamily\small
\noindent
function powmod(a, n, m):\\
\quad result := 1\\
\quad a := a mod m\\
\quad while n > 0:\\
\qquad if n is odd:\\
\quad\qquad result := (result $\ast$ a) mod m\\
\qquad a := (a $\ast$ a) mod m\\
\qquad n := n div 2\\
\quad return result
\end{minipage}
\end{center}

\begin{example}{}{ex:p1-5}
Вычислим $7^{13} \bmod 23$. Запишем $13 = 1101_2$, то есть $13 = 8 + 4 + 1$. Значит, $7^{13} = 7^{8}\cdot 7^{4}\cdot 7^{1}$. Считаем последовательно по модулю 23:
\[
7^1 = 7,\quad 7^2 = 49 \equiv 3,\quad 7^4 \equiv 3^2 = 9,\quad 7^8 \equiv 9^2 = 81 \equiv 81 - 3\cdot 23 = 12.
\]
Тогда $7^{13} \equiv 12\cdot 9\cdot 7 = 756 \equiv 756 - 32\cdot 23 = 756 - 736 = 20 \modn{23}$. Всего --- около пяти умножений вместо тринадцати.
\end{example}

\paragraph*{Та же идея в другом месте: числа Фибоначчи.} Знакомая со школы последовательность Фибоначчи задаётся рекуррентно:
\[
F_0 = 0,\quad F_1 = 1,\quad F_{n+1} = F_n + F_{n-1}.
\]

На первый взгляд, чтобы вычислить $F_n$, нужно сделать $n$ сложений --- линейное время. Однако и здесь работает «разделяй и властвуй»! Запишем рекуррентное соотношение в матричной форме:
\[
\begin{pmatrix} F_{n+1} \\ F_n \end{pmatrix}
\;=\;
\begin{pmatrix} 1 & 1 \\ 1 & 0 \end{pmatrix}
\begin{pmatrix} F_n \\ F_{n-1} \end{pmatrix}.
\]
Обозначим матрицу как $A$. Применяя её $n$ раз к вектору $(F_1, F_0)^\top = (1,0)^\top$, получим:
\[
\begin{pmatrix} F_{n+1} \\ F_n \end{pmatrix}
\;=\;
A^n
\begin{pmatrix} 1 \\ 0 \end{pmatrix}.
\]
Прямое доказательство --- индукцией по $n$.

\begin{importantbox}
Тем самым задача сводится к возведению матрицы $2{\times}2$ в степень $n$. Но матрицы можно перемножать «делением пополам» точно так же, как числа! Сложность --- $O(\log n)$ матричных умножений, то есть около $\log_2 n$ операций; вместо линейного числа шагов получается логарифмическое (рис.~\ref{fig:p1-3}).
\end{importantbox}

\begin{figure}[H]
\centering
\begin{tikzpicture}[scale=1.0, every node/.style={font=\small}]
  % Сравнение линейного и логарифмического подхода
  \begin{scope}[xshift=-5cm]
    \node[font=\bfseries\color{prosvRed}] at (1.8,2.4) {Наивно: $n$ шагов};
    \foreach \i in {0,...,9}{
      \fill[prosvRed!60!white] (\i*0.35,0) rectangle (\i*0.35+0.3,0.5);
    }
    \node at (1.5, -0.4) {$F_1 \to F_2 \to \cdots \to F_{10}$};
  \end{scope}
  \begin{scope}[xshift=1cm]
    \node[font=\bfseries\color{prosvGreen}] at (1.8,2.4) {С «делением пополам»: $\log_2 n$ шагов};
    \foreach \i in {0,...,3}{
      \fill[prosvGreen!70!white] (\i*0.5,0) rectangle (\i*0.5+0.45,0.5);
    }
    \node at (1.5,-0.4) {всего $\approx \log_2 10 \approx 4$ шага};
  \end{scope}
\end{tikzpicture}
\caption{Сравнение наивного и «разделяй и властвуй»-подхода к вычислению $F_n$.}
\label{fig:p1-3}
\end{figure}

Этот пример замечателен тем, что показывает: даже линейная на первый взгляд задача может «прятать» внутри себя структуру, ускоряющую решение экспоненциально.

% -----------------------------------------------------------------
\subsection{Функция Эйлера и малая теорема Ферма}

Чем дальше, тем интересней. Сейчас мы введём одну функцию, играющую ключевую роль в современной криптографии.

\begin{definition}{}{def:p1-5}
\term{Функцией Эйлера} $\Eul(n)$ называется количество натуральных чисел в отрезке $[1, n]$, взаимно простых с $n$.
\end{definition}

\begin{example}{}{ex:p1-6}
$\Eul(1) = 1$, $\Eul(2) = 1$, $\Eul(6) = 2$ (а именно, числа 1 и 5), $\Eul(10) = 4$ (числа 1, 3, 7, 9).
\end{example}

\paragraph*{Случай простого числа.} Если $p$ --- простое число, то с $p$ \emph{не} взаимно прост лишь сам $p$ среди чисел $1, 2, \ldots, p$. Значит,
\[
\boxed{\;\Eul(p) = p - 1\;}\quad \text{для простого } p.
\]

\paragraph*{Случай произведения двух простых.} Пусть $n = p\cdot q$, где $p, q$ --- различные простые. Какие числа из $[1, n]$ не взаимно просты с $n$? Только кратные $p$ или кратные $q$. Кратных $p$ ровно $q$ штук (это $p, 2p, \ldots, qp$); кратных $q$ ровно $p$ штук; единственное число, кратное обоим, --- $pq = n$. По формуле включений-исключений:
\[
\Eul(pq) = pq - p - q + 1 = (p-1)(q-1).
\]
Это короткое наблюдение --- сердце алгоритма RSA, который мы изучим в §~\ref{sec:rsa-dh}.

\begin{example}{}{ex:p1-7}
$\Eul(15) = \Eul(3\cdot 5) = 2\cdot 4 = 8$. Действительно, взаимно простыми с 15 в отрезке $[1,15]$ являются числа $\{1, 2, 4, 7, 8, 11, 13, 14\}$ --- ровно 8 чисел.
\end{example}

В общем виде: для $n = p_1^{\alpha_1}\cdots p_k^{\alpha_k}$ выполняется $\Eul(n) = n\prod_{i=1}^k\!\bigl(1 - \tfrac{1}{p_i}\bigr)$, но эта формула нам сегодня понадобится только в указанных двух частных случаях.

\paragraph*{Малая теорема Ферма.} Один из самых красивых результатов классической теории чисел.

\begin{theorem}{малая теорема Ферма, 1640 г.}{thm:p1-3}
\label{thm:fermat}
Пусть $p$ --- простое число и $a$ --- целое, не делящееся на $p$. Тогда
\[
a^{p-1} \equiv 1 \modn p.
\]
Равносильно: для любого целого $a$ выполняется $a^p \equiv a \modn p$.
\end{theorem}

\begin{proof}[Схема доказательства]
Рассмотрим набор $\{a, 2a, 3a, \ldots, (p-1)a\}$ --- $(p-1)$ ненулевых остатков, умноженных на $a$. Все они различны по модулю $p$: если $ia \equiv ja \modn p$, то $a(i-j) \equiv 0 \modn p$, а так как $a$ не делится на $p$ и $|i-j| < p$, получаем $i = j$.

Значит, набор $\{a, 2a, \ldots, (p-1)a\}$ --- это в точности (с точностью до порядка) набор $\{1, 2, \ldots, p-1\}$ по модулю $p$. Перемножим всё:
\[
a\cdot 2a\cdot 3a \cdots (p-1)a \;\equiv\; 1\cdot 2\cdot 3\cdots(p-1) \modn p,
\]
то есть $a^{p-1}\cdot (p-1)! \equiv (p-1)! \modn p$. Поскольку $(p-1)!$ взаимно прост с $p$, его можно «сократить» --- получим $a^{p-1} \equiv 1 \modn p$. \qedhere
\end{proof}

Обобщение --- \term{теорема Эйлера}: $a^{\Eul(n)} \equiv 1 \modn n$ для любого $a$, взаимно простого с $n$. Доказывается аналогично, только вместо чисел $1, \ldots, p-1$ берутся все числа из $[1,n]$, взаимно простые с $n$.

\begin{interestbox}
Пьер Ферма, по профессии юрист, посвящал математике вечера. В 1640 году он написал письмо своему другу с этим утверждением, заметив: «Я бы прислал доказательство, если бы не боялся, что оно слишком длинное». Первое известное полное доказательство опубликовал Леонард Эйлер --- спустя сто лет.
\end{interestbox}

% -----------------------------------------------------------------
\subsection{Сколько вокруг простых чисел?}

Прежде чем переходить к криптографии, ответим на важный практический вопрос: \emph{часто ли встречаются простые числа?} Если простых чисел мало --- ими нельзя будет полноценно пользоваться.

\begin{theorem}{Чебышёв -- Адамар -- Валле-Пуссен, асимптотический закон распределения простых чисел}{thm:p1-4}
Пусть $\pi(N)$ --- количество простых чисел, не превосходящих $N$. Тогда
\[
\pi(N) \;\sim\; \frac{N}{\ln N} \qquad \text{при } N \to \infty.
\]
\end{theorem}

\noindent
В пересчёте на «вероятность»: если случайно взять натуральное число вокруг $N$, оно окажется простым примерно с шансом $1/\ln N$. Например, среди 1000-значных чисел простыми оказывается примерно одно из $\ln 10^{1000} = 1000\ln 10 \approx 2300$. Это совсем немало (рис.~\ref{fig:p1-4}).

\begin{figure}[H]
\centering
\begin{tikzpicture}
\begin{axis}[
  width=11cm, height=6cm,
  xlabel={$N$}, ylabel={число простых $\le N$},
  xmin=1, xmax=1000, ymin=0, ymax=200,
  axis lines=left,
  grid=major, grid style={prosvGray!30},
  legend pos=north west, legend cell align=left,
  xlabel style={font=\small}, ylabel style={font=\small},
  tick label style={font=\scriptsize},
  every axis plot/.append style={very thick}
]
% Точная функция pi(N) приближённо
\addplot[prosvRed, domain=2:1000, samples=80] {x/ln(x)};
\addlegendentry{$\dfrac{N}{\ln N}$ (приближение)};
% Фактические точки pi(N) для N=10,50,100,...
\addplot+[
  only marks, mark=*, prosvBlue, mark size=1.4pt
] coordinates {
  (10,4) (50,15) (100,25) (200,46) (300,62) (400,78) (500,95)
  (600,109) (700,125) (800,139) (900,154) (1000,168)
};
\addlegendentry{точное значение $\pi(N)$};
\end{axis}
\end{tikzpicture}
\caption{Закон распределения простых чисел: количество простых чисел до $N$ ведёт себя как $N/\ln N$.}
\label{fig:p1-4}
\end{figure}

Практический смысл: чтобы найти простое число длиной, скажем, 1024 бита для RSA-ключа, нужно перебрать в среднем порядка 700 случайных чисел и каждое проверить на простоту. И это нас приводит к следующему вопросу.

% -----------------------------------------------------------------
\subsection{Проверка простоты и задача факторизации}
\label{subsec:primality}

Два «двойственных» вопроса:
\begin{description}
  \item[Проверка простоты:] дано число $n$ --- простое оно или составное?
  \item[Факторизация:] дано составное $n$ --- разложить его на простые множители.
\end{description}

Может показаться, что эти задачи примерно одинаковой сложности. Но это не так: \emph{они принципиально разной сложности}, и именно на этом контрасте основана вся современная криптография.

\paragraph*{Тест Ферма.} Если $p$ простое, то по теореме~\ref{thm:fermat} для любого $a$, не кратного $p$, выполняется $a^{p-1} \equiv 1 \modn p$. Возьмём это как \emph{пробный признак}: если для случайного $a$ выполнено $a^{n-1} \not\equiv 1 \modn n$, то $n$ заведомо \emph{не простое}. Если же $a^{n-1} \equiv 1 \modn n$ --- то $n$ простое \emph{с большой вероятностью} (но иногда бывают и составные числа, проходящие тест --- так называемые \term{числа Кармайкла}).

\paragraph*{Тест Миллера--Рабина.} Усовершенствованный вариант теста Ферма. Он использует ту же идею, но дополнительно «отсеивает» числа Кармайкла. Принципиальные свойства:
\begin{itemize}
  \item это \term{вероятностный} тест: при отрицательном ответе число точно составное; при положительном --- простое с вероятностью ошибки не больше $1/4$ за один проход;
  \item делая $k$ независимых проходов, мы уменьшаем вероятность ошибки до $4^{-k}$. Уже при $k = 20$ это меньше, чем $10^{-12}$;
  \item работает за время $O(k\cdot \log^3 n)$ --- то есть \term{полиномиально} от числа цифр.
\end{itemize}

\paragraph*{Тест AKS (2002 г.).} Манидра Агравал, Нираж Каял и Нитин Саксена --- индийские математики --- впервые предъявили \emph{детерминированный} полиномиальный алгоритм проверки простоты. На практике он работает медленнее Миллера--Рабина, но важен как теоретическое достижение: показал, что класс задач, разрешимых за полиномиальное время, включает в себя проверку простоты.

\begin{interestbox}
Открытие AKS-алгоритма стало сенсацией: один из руководителей --- Нитин Саксена --- был на тот момент аспирантом первого года. Результат опубликован в 2004 году в одном из самых престижных математических журналов --- «Annals of Mathematics» --- и был назван «алгоритмом тысячелетия» в области теории чисел.
\end{interestbox}

\paragraph*{А что с факторизацией?} Тут картина совсем другая. Лучший известный сегодня алгоритм --- \emph{общий метод решета числового поля} --- работает за время, растущее как
\[
\exp\!\bigl(c\cdot \sqrt[3]{(\ln n)(\ln\ln n)^{2}}\bigr),
\]
что по сути \term{субэкспоненциально}, но \emph{не полиномиально}. Для 1024-битных $n$ это означает миллиарды лет работы суперкомпьютера.

\begin{importantbox}
Колоссальный разрыв в сложности: \emph{проверить, простое ли число} --- легко (полиномиально), а \emph{разложить число на простые множители} --- невероятно трудно (субэкспоненциально). Именно этот разрыв и эксплуатируется в RSA и других криптосистемах. Безопасность интернета прямо опирается на то, что у человечества пока нет быстрого алгоритма факторизации.
\end{importantbox}

\paragraph*{А если появится квантовый компьютер?} В 1994 году Питер Шор предложил квантовый алгоритм, факторизующий числа за полиномиальное время. Реальных квантовых компьютеров достаточного масштаба пока нет (на 2025 год удавалось факторизовать лишь крошечные числа в качестве демонстрации). Но если такие компьютеры будут построены --- современный RSA рухнет, и потребуются принципиально новые криптографические схемы (это направление называется \engterm{постквантовая криптография}).

% -----------------------------------------------------------------
\subsection*{Что мы получили}

Подведём итог. На вопрос «зачем нам всё это в учебнике информатики» мы готовы ответить: только что мы построили \emph{инструменты}, без которых не работает ни один защищённый сайт. А именно:
\begin{itemize}
  \item \term{сравнения по модулю} --- язык, на котором будем описывать шифрование;
  \item \term{расширенный алгоритм Евклида} --- быстро ищет обратный элемент по модулю;
  \item \term{быстрое возведение в степень} --- позволяет работать с гигантскими степенями;
  \item \term{малая теорема Ферма} и \term{функция Эйлера} --- ключ к схеме RSA;
  \item \term{разрыв в сложности} между проверкой простоты и факторизацией --- источник «безопасности».
\end{itemize}

\noindent
В следующих параграфах мы соберём из этих кирпичиков работающие криптографические протоколы.

% -----------------------------------------------------------------
\begin{tasksbox}
\begin{enumerate}
  \item Найдите $\NOD(231, 165)$ алгоритмом Евклида. Представьте НОД в виде $231 x + 165 y$.
  \item Найдите $5^{-1} \modn{37}$. Используйте расширенный алгоритм Евклида.
  \item Вычислите $3^{100} \bmod 7$, пользуясь малой теоремой Ферма.
  \item Сколько существует чисел в отрезке $[1, 100]$, взаимно простых с $100$? (Подсказка: $100 = 2^2\cdot 5^2$, но вы можете использовать формулу для $\Eul(p^2)$ или прямой подсчёт.)
  \item Запишите быстрое возведение в степень в виде программы на знакомом вам языке программирования.
  \item* Докажите, что для составного $n$ всегда найдётся $a$ из $[2, n-1]$, для которого $a^{n-1} \not\equiv 1 \modn n$ --- кроме чисел Кармайкла. Найдите наименьшее число Кармайкла.
  \item* С помощью матричной формулы для чисел Фибоначчи вычислите $F_{20}$.
\end{enumerate}
\end{tasksbox}
